% Generated draft: requires mathematical and source review.
\documentclass[11pt]{article}
\usepackage[margin=25mm]{geometry}
\usepackage{fontspec}
\setmainfont{Latin Modern Roman}
\usepackage{amsmath,amssymb,mathtools,unicode-math}
\setmathfont{Latin Modern Math}
\usepackage{enumitem,needspace,xurl,hyperref}
\setlength{\parindent}{0pt}
\setlength{\parskip}{5pt}
\setlength{\emergencystretch}{3em}
\newcommand{\N}{\mathbb{N}}
\newcommand{\Z}{\mathbb{Z}}
\newcommand{\Q}{\mathbb{Q}}
\newcommand{\R}{\mathbb{R}}
\newcommand{\C}{\mathbb{C}}
\newcommand{\F}{\mathbb{F}}
\newcommand{\A}{\mathbb{A}}
\newcommand{\PP}{\mathbb{P}}
\newcommand{\E}{\mathbb{E}}
\newcommand{\eps}{\varepsilon}
\newcommand{\dd}{\,\mathrm{d}}
\begin{document}
\section*{UnsolvedMath ID: 30005570}
\textbf{Source problem code:} catalogue entry
\subsection{Definitions and mathematical statement}
Let \(f\) be a fixed modular form of weight \(k\geq 2\), and let \(d\geq 1\) be a fixed integer, as in the supplied statement. A Dirichlet character modulo \(q\) is a character of \((\Z/q\Z)^\times\), extended periodically to integers coprime to \(q\) and by zero to the other integers. Its conductor \(\operatorname{cond}(\chi)\) is the least modulus from which it is induced. It is primitive when its modulus equals its conductor. Its order \(\operatorname{ord}(\chi)\) is its exact multiplicative order on the unit group.

For real \(X\geq 1\), define
\[
K_d(X)=\left\{\chi:\ \chi\text{ is a primitive Dirichlet character},\quad \operatorname{ord}(\chi)=d,\quad \operatorname{cond}(\chi)\leq X\right\}.
\]
Each primitive character is counted at its conductor, without also counting its induced realizations at larger moduli. Exact order is intended, rather than order dividing \(d\). The range \(d\geq 1\) includes the trivial primitive character of conductor \(1\).

Write \(L(f,\chi,s)\) for the Dirichlet twist of the \(L\)-function associated with \(f\), with the normalization intended by the source placing the central argument at \(s=k/2\). Where \(f\) has Fourier coefficients \(a_f(n)\), the usual arithmetic convention starts from
\[
L(f,\chi,s)=\sum_{n\geq 1}a_f(n)\chi(n)n^{-s}
\]
in a half-plane of convergence, followed by analytic continuation. The supplied record does not specify the class of modular forms, level, character, coefficient normalization, or analytic hypotheses needed to make every stated central value a finite, well-defined value. These are unresolved qualifications; no cuspidality, eigenform, newform, or nonzero hypothesis is added here.

\textbf{Requested asymptotic.} For every fixed pair \((f,d)\) covered by the intended hypotheses, determine the asymptotic behavior, as \(X\to\infty\), of
\[
N_{f,d}(X)=\#\left\{\chi\in K_d(X):L(f,\chi,k/2)\neq 0\right\}.
\]
This formulation presupposes that the displayed central values are defined and finite. The source supplies neither a proposed main term nor a universally asserted limiting proportion.

The original extract also discusses the elliptic-curve specialization \(L(f,s)=L(E,s)\), with central argument \(1\). It describes Goldfeld's conjecture as predicting that, among quadratic characters ordered by conductor,
\[
\operatorname{ord}_{s=1}L(f,\chi,s)=0\quad\text{for }50\%,\qquad
\operatorname{ord}_{s=1}L(f,\chi,s)=1\quad\text{for }50\%.
\]
Here the order of vanishing is the multiplicity of the zero of the analytically continued function at \(s=1\). This is supplied conjectural background for the quadratic elliptic-curve setting, not an asserted answer for arbitrary \(f\) and \(d\). The extract is truncated around this specialization, and its reference \([3]\) is not identified in the record.
\subsection{Short English statement}
For a fixed modular form \(f\) of weight \(k\geq 2\) and a fixed positive integer \(d\), count primitive Dirichlet characters of exact order \(d\), ordered by conductor, whose twists have nonzero central value \(L(f,\chi,k/2)\). Determine how this count grows as the conductor bound tends to infinity. The precise hypotheses on \(f\) and the analytic meaning of these central values require clarification. The supplied Goldfeld discussion concerns the quadratic elliptic-curve specialization.
\subsection{Sources}
\begin{itemize}
\item Local UnsolvedMath registry: \texttt{data/unsolvedmath/problems.json}, permanent ID \(30005570\), title ``Nonvanishing Statistics of Twisted Central L-Values,'' source code \texttt{catalogue entry}; collection: Oberwolfach Reports Open Problems.
\item Supplied report citation: \emph{Algebraische Zahlentheorie} (2023), Oberwolfach Reports, pp.\ 1543--1606, DOI \texttt{10.4171/owr/2023/28}. \url{https://doi.org/10.4171/owr/2023/28}. The record attributes the relevant contribution, ``Horizontal p-adic L-functions with applications to L-values,'' to Daniel Kriz, reporting joint work with Asbjørn Nordentoft. No separate proposer is named.
\item Supplied literature source: Kriz--Nordentoft, \emph{Horizontal p-adic L-functions}. \url{https://arxiv.org/abs/2310.20678}.
\item Supplied literature source: Kriz, \emph{Horizontal p-adic L-functions with applications to L-values}. \url{https://doi.org/10.4171/owr/2023/28}.
\item Literature assessment dated exactly \texttt{2026-08-22}: the supplied curation labels the problem open, states that the general nonvanishing asymptotic is unknown and Goldfeld's prediction remains conjectural, and reports strong lower bounds and special-case progress from horizontal p-adic methods. This is supplied dated curation, not an independently verified literature assessment.
\item The registry labels the statement \texttt{corrected\_verified} and says it was checked against the official report. That verification is a supplied provenance claim, not verification performed for this draft. The original extract cites Goldfeld's conjecture as reference \([3]\), but supplies no bibliographic entry for that label.
\end{itemize}
\end{document}
